% !TeX root = ../组合数学笔记.tex
% 由主文档通过 \input 载入；请编译主文档。
\setcounter{chapter}{16} % 保留原书章号。
\chapter{组合算法与综合问题}
\label{ch:algorithms}

计数公式、递推和存在性证明可以导出算法，但必须说明算法实际输出什么。输出一个对象、对象总数或全部对象，其资源需求可能相差很大。本章默认输入为有限离散数据，先按算术操作数分析成本，再说明整数位长与输出规模的影响。

\section{枚举、排名与反排名}

枚举长度为 $n$ 的零一串，可以逐位选择 $0$ 或 $1$；枚举排列则逐位选择一个尚未使用的标签。每个合法对象都有唯一的选择路径，所以无重复也无遗漏；每次增加已填位置数，深度至多为 $n$，故过程终止。输出全部 $n!$ 个排列、每个明确写出 $n$ 个元素时，仅写入输出就需要至少与 $n\cdot n!$ 同阶的操作，不能期待关于 $n$ 的多项式总时间。

固定从零开始的字典序编号，\emph{排名（ranking）}将对象映射到其编号，\emph{反排名（unranking）}从编号恢复对象。对排列 $\pi=(\pi_1,\ldots,\pi_n)$，令 $c_i$ 为尚未使用且小于 $\pi_i$ 的标签数，则
\[
 \operatorname{rank}(\pi)=\sum_{i=1}^n c_i(n-i)!.
\]
因为在第 $i$ 位第一次选择一个较小标签时，后面有 $(n-i)!$ 种补全，且不同首次分歧互斥。逆向给定 $0\le R<n!$，在第 $i$ 步取剩余有序标签中从零编号为 $\lfloor R/(n-i)!\rfloor$ 的元素，再令 $R$ 为相应余数，即可恢复排列。每次商都在合法范围内，商余分解也说明二者互逆。例如 $3142$ 的 $c_i$ 为 $2,0,1,0$，排名为 $13$。

\section{动态规划与状态压缩}

\emph{动态规划（dynamic programming）}将重复出现的子问题只计算一次。设不同正整数面值为 $a_1,\ldots,a_k$，目标金额为 $T\ge0$，各面值不限枚数、同面值硬币不可区分。令 $D_j(t)$ 计数只用前 $j$ 种面值凑成 $t$ 的方案，规定 $D_0(0)=1,D_0(t)=0$（$t>0$），并将负金额项置零。按是否使用第 $j$ 种硬币分类，得到
\[
 D_j(t)=D_{j-1}(t)+D_j(t-a_j).
\]
后一项删除一枚第 $j$ 种硬币后仍允许使用该种，因此是 $D_j$ 而非 $D_{j-1}$。一维实现如下：
\begin{verbatim}
CountChange(a[1..k], T):
    dp[0..T] = 0
    dp[0] = 1
    for j = 1..k:
        for t = a[j]..T in increasing order:
            dp[t] = dp[t] + dp[t-a[j]]
    return dp[T]
\end{verbatim}
外层循环不变式是：处理完第 $j$ 种面值，数组为 $D_j$；内层递增保证新项使用当前层的较小金额。算法用 $O(kT)$ 次加法和 $O(T)$ 个整数存储；若 $T$ 以二进制给出，这对输入位长不一定是多项式时间。把金额作为外循环、面值作为内循环则按最后一枚分类，计数的是有序投币串；例如面值 $1,2$ 凑成 $3$，无序方案为两种，有序串为三种。

对 $n$ 人与 $n$ 个不同任务的受限分配，令第 $i$ 人可选任务集为 $A_i\subseteq[n]$，状态 $F(S)$ 表示把前 $|S|$ 人分配到任务集 $S$ 的方案数。初值为 $F(\varnothing)=1$，其余初始为零；按集合大小递增，对每个 $j\in A_{|S|+1}\setminus S$ 更新 $F(S\cup\{j\})\mathrel{+}=F(S)$。最后输出 $F([n])$。删除最后一人的任务唯一确定前驱，所以转移不重不漏；有 $2^n$ 个状态，每态至多 $n$ 条转移，需 $O(n2^n)$ 次加法和 $O(2^n)$ 个整数。答案至多为 $n!$，单个数可能需要 $O(n\log n)$ 位，这一位成本应与操作数分开。

\section{回溯、剪枝与对称消除}

\emph{回溯（backtracking）}沿部分解的搜索树逐步扩展；所有合法分支都被访问时保证完备。剪枝必须有“当前前缀没有合法补全”的证明。例如受限分配的剩余二分图若不能饱和剩余人员，则整个分支不可能成功；可以用增广路算法判定，而仅观察每个人都有候选任务是不够的。

若只需对称等价类，可为完整对象建立编码，并保留其轨道中编码字典序最小者。有限轨道总有唯一最小编码，所以每个等价类保留一次；代价是对每个对象检查变换。更高效的部分对称消除必须证明每个轨道仍保留适当数量的代表，不能随意固定某个位置后就把全部计数机械乘回群大小。

直接枚举适合核对小参数，也能发现漏项、边界和重数错误；但有限范围内全部吻合，仍不能证明一般公式。反之，若枚举结果不同，应先核对双方的对象、标签、顺序和对称约定，再判断是哪一个推导出错。

\section{综合专题与方法选择}

对 $x_1+x_2+x_3=8$、$0\le x_i\le4$，容斥给出 $45-3\cdot10=15$；生成函数给出 $[x^8](1+x+\cdots+x^4)^3$；动态规划令 $F_j(t)=\sum_{u=0}^4F_{j-1}(t-u)$，初值 $F_0(0)=1$，同样得 $F_3(8)=15$。容斥适合坏事件交集容易计数时，生成函数适合独立数量限制，动态规划则直接给出可计算的整数表。

六个圆周位置恰有两个红色，其余蓝色，识别旋转后可以按两红点沿圆周的较短距离 $1,2,3$ 分类，恰有三类。Burnside 计算中，恒等旋转固定 $\binom62=15$ 个图案，半周旋转固定三种相对红点，其余旋转不可能保持恰好两红，故平均为 $(15+3)/6=3$。枚举所有十五个二元子集并取旋转规范代表，可以作为独立核对。

对三人的候选集 $\{a,b\},\{a,b\},\{b,c\}$，Hall 定理回答可行性，最大匹配给出一个分配，子集动态规划则给出方案总数为二：第三人必须取 $c$，前两人交换 $a,b$。最大匹配大小本身并不等于匹配方案数。先确定问题输出，才能选对算法。

\section{习题}
\begin{enumerate}
  \item\label{ex:algorithms-1} 计算排列 $2413$ 的零起点字典序排名，并从编号 $13$ 反排名恢复四元素排列。
  \item\label{ex:algorithms-2} 写出不含连续三个 $1$ 的零一串计数算法，给出状态、初值、循环顺序、终止与复杂度。
  \item\label{ex:algorithms-3} 用 $A_1=\{1,2\},A_2=\{2,3\},A_3=\{1,3\}$ 建立子集动态规划，求完整分配数。
  \item\label{ex:algorithms-4} 选择一个小规模计数问题，分别用独立枚举和一种公式方法求解；说明二者相符验证了哪些内容，仍未证明哪些内容。
\end{enumerate}

